\section{Retained complete-record theorem and its proof route}
\label{app:v58-retained-front}


\paragraph{Record, source and arithmetic.}
Let $T_R$ be the collision map of the triangular Lorentz table,
$R\in[0.45,0.47]$, and $F_R^*$ its actual first return to the specified
section $Y_R^*$. Put $c=\nu(Y_R^*)=91/(10000\pi)$ and
$\nu_R^*=c^{-1}\nu|_{Y_R^*}$. At a fixed return index $n$ the original
record $J_{n,R}=(K_{n,R},N_{n,R},T_{n,R})$ has the complete density
\[
 \nu_R^*\{K_{n,R}=k,N_{n,R}=m,T_{n,R}\in du\}
                       =p_{n,R}(k,m,u)\,du.
\]
The measure is counting in $(k,m)\in\Z^2\times\{n,n+1,\ldots\}$
and Lebesgue in $u>0$. The finite Gaussian peak sum
$\mathcal L_{m,R}(k_1,k_2,u,n)$ is defined in
\eqref{eq:v41-transition-kernel}. It retains every occupation residue
and every spectrally damped branch needed for parameter-uniformity.
Define $\gamma_{n,R}=m^{-2}\mathcal L_{m,R}$ and let
$\mathsf A_{n,R}$ be its normalized positive part on this whole mixed
space, as in \eqref{eq:v57-full-reference}.

\clearpage
\begin{leadtheorem}[The full arithmetic record and its coupled bridges]
\label{thm:intro-v57-global}
Uniformly in $R$, as the actual return index $n\to\infty$,
\[
 \sum_{m=n}^\infty\sum_{k\in\Z^2}\int_0^\infty
   |p_{n,R}(k,m,u)-m^{-2}\mathcal L_{m,R}(k_1,k_2,u,n)|\,du\to0.
\]
The total mass of $(\gamma_{n,R})_+$ tends to one. Consequently
$d_{\rm TV}(\operatorname{Law}_{\nu_R^*}(J_{n,R}),\mathsf A_{n,R})\to0$.
Here probability total variation is one half of variation mass.
No average of return indices is taken.

Let $\mathcal B_{N_{n,R},R}$ and $\mathcal Y_{n,R}$ be the collision
and actual-return paths pinned to their own terminal record. Put
$A_R=c^{-1/2}L_R^{-1}$, where $L_R$ is the covariance change in
\eqref{eq:v37-covariance-change}, and set
\[
 \mathsf H_R=\operatorname{Law}(\mathbb B_{\Omega_R},
                                      A_R\mathbb B_{\Omega_R}).
\]
Then
\[
 \left\|\operatorname{Law}_{\nu_R^*}
  (J_{n,R},\mathcal B_{N_{n,R},R},\mathcal Y_{n,R})
                 -\mathsf A_{n,R}\otimes\mathsf H_R
       \right\|_{\mathrm{obs},\mathrm{BL}}\to0
\]
uniformly. The norm allows every bounded measurable test of the record,
with a unit bounded-Lipschitz test of the two paths chosen measurably
from that record. It is not total variation in the path coordinates.
The two Gaussian bridges have covariances $\Omega_R,D_R$ and are
coupled by $A_R$, not independently sampled.

If $\Delta_n$ denotes the uniform error in the last display, applying
the same arbitrary observation kernel to both records cannot increase
it. Conditioning on the full output event of reference probability at
least $r_n>\Delta_n$ gives output total-variation error at most
$\Delta_n/(r_n-\Delta_n)$ and joint path-dual error at most
$2\Delta_n/(r_n-\Delta_n)$. In particular these errors vanish when
$\Delta_n/r_n\to0$. No polynomial return-count rate is asserted.
\end{leadtheorem}
\begin{proof}
Theorems~\ref{thm:v57-global-raw-TV}, \ref{thm:v57-global-coupled}
and~\ref{thm:v57-selection} prove the assertions.
\end{proof}

\paragraph{The new spectral estimate.}
For the retained peak data in \eqref{eq:v41-peak-data}, we prove
\[
 |\mu_j(R)-\mu_R|^2\le C\kappa_j(R),\qquad
 |\mathcal L_{m,R}(y)|\le C e^{-a|y-m\mu_R|^2/m}.
\]
The first inequality is stronger than continuity at a resonance. Its
proof dominates a branch with a nonzero smooth physical pairing by the
positive exponential moment at an imaginary tilt, then optimizes that
tilt. It does not require a nonzero section residue. The second
inequality prevents weakly damped peaks from carrying mass to remote
centers as the radius varies. It justifies summation of all collision
and displacement labels and integration over the full roof line.
The original return fourth moment controls the source tail.

\paragraph{Arithmetic is part of the theorem.}
At each fixed radius the normalized reference may equivalently use
\[
 n^{-2}\mathfrak a_R(k,n,m)g_{D_R}
  \left(\frac{k_1,k_2,m-n/c,u-n\bar\tau_R/c}{\sqrt n}\right),
\]
by Corollary~\ref{cor:v57-fixed-radius-raw}. Zero residue classes are
retained. The uniform-in-radius theorem uses the full transition
kernel instead of assuming that arithmetic types are constant.
An unmodulated Gaussian specialization is not asserted without
verification of the concrete section residue criterion.

\paragraph{The unrestricted pointwise target.}
The raw object remains this same four-coordinate return record. The
positive identity of Theorem~\ref{thm:v51-positive-error} reduces its
unrestricted two-sided pointwise theorem to the essential heights of
both the incidence and clearance sources at the original exact labels.
Those requirements are not consequences of the total-variation result
above. The local weak $L^{145/144}$ law, positive lower law and common
all-resolution exceptional sets remain stated and proved, together
with their exact quantifiers, in
Appendix~\ref{app:v57-retained-front} and the cited core sections.
Every inherited mathematical module remains compiled.

\paragraph{Proof route and comparison.}
Sections~\ref{sec:v57-transition-tails} and~\ref{sec:v57-coupled-bridges}
give the new route first. They explicitly identify the inherited
occupation spectrum, local raw variation and joint clock identity used
at each step. The later parts retain the complete geometric and raw
inversion proofs. The source analysis goes beyond a fixed interval
limit because its roof tests are allowed to depend measurably on all
exact labels. The new passage from local to global needs an arithmetic
reference-tail theorem, not merely tightness of the true process.
Positive-pressure domination, Gaussian completion of the square and
postselection normalization are not claimed as new general methods.
Their use here is to keep moving occupation peaks, the complete raw
record and the two coupled clocks in one statement. The theorem-level
comparisons with \cite{DPZ,DN} remain in
Section~\ref{sec:v35-literature-routes}; neither cited interval theorem
is differentiated to obtain a density value.

